\section{Methodology}

\subsection{Scorer (Inherited)}

\begin{equation}
\begin{aligned}
S(h,c) =\;& \alpha\,\mathrm{EPSS}(c) + \beta\,\mathrm{HRS}(h) \\
          & + \gamma\,\mathrm{KEV\_recency}(c) \\
          & + \delta\,\bigl(\mathrm{EPSS}(c) \cdot \mathrm{HRS}(h)\bigr).
\end{aligned}
\label{eq:hp}
\end{equation}

\subsection{CUSUM Detector}

At each window $w \geq 2$:
\begin{enumerate}
\item Compute the calibration-target shift $\Delta_w$ as in
Paper~10~\cite{paper10adaptive}.
\item Update the accumulator:
\begin{equation}
C_w = \max(0,\; C_{w-1} + |\Delta_w| - k), \quad C_1 = 0.
\end{equation}
\item Fire if $C_w \geq h$.
\item On fire: revert \textsf{cusum} to Paper~4 fixed weights for
window $w$, then reset $C_w = 0$.
\item Otherwise: use the lag-1 fitted weights at window $w$.
\end{enumerate}

Pre-registered $k = 0.04$, $h = 0.10$. The slack $k$ is one-half the
typical $|\Delta_w|$ at K=50 from Paper~11's published calibration
log; the threshold $h$ is the largest single-$\tau$ value from
Paper~11's grid.

\subsection{Why $C_w$ Resets to Zero on Fire}

Reset is the standard one-sided CUSUM convention~\cite{page1954cusum}.
After firing, the detector treats the post-fire window as a fresh
operating regime and only re-fires if a new cumulative excursion
crosses $h$. Without reset the detector would fire repeatedly on the
same persistent shift, which is the behaviour the magnitude detector
of Papers~10--12 already provides.

\subsection{Other Strategies (Inherited)}

\textsf{lag1}, \textsf{offline}, and \textsf{gated} are unchanged
from Papers~7--12. \textsf{cap\_aware} uses the Paper~12
$\tau_K = \{50: 0.20, 100: 0.05, 200: 0.02\}$ vector and is
included as a directly-comparable control.
